\chapter{Background and State of the Art}
\label{chap:background}

This chapter develops the concepts the thesis is built from, as a
progression: the distributed-systems fundamentals that define the
problem's vocabulary (\S\ref{sec:bg-ds}); distributed tracing, the discipline
that first made cross-machine causality observable
(\S\ref{sec:bg-tracing}); the distributed shared log, the storage abstraction
this thesis adopts as its substrate (\S\ref{sec:bg-log}); ChronoLog, the
concrete log store used (\S\ref{sec:bg-chronolog}); LLM agents and
multi-agent systems, the workload being observed (\S\ref{sec:bg-agents});
and finally the state of the art, whose two disjoint camps define the gap
this thesis targets (\S\ref{sec:bg-sota}). Most sections close by noting
what the thesis takes from them.

\section{Distributed Systems Fundamentals}
\label{sec:bg-ds}

\subsection{Time, Ordering, and Causality}
\label{sec:bg-time}

A distributed system is a set of processes on separate machines that
communicate only by exchanging messages. Its defining difficulty is that
there is no global ``now'': each machine has its own clock, clocks drift, and
no observer sees the whole system at once. Any statement about what the
system \emph{is doing} must therefore be reconstructed from locally
recorded observations.

Lamport's insight is that such reconstruction turns not on wall-clock
simultaneity but on \emph{causality}, captured by the
\emph{happened-before} relation~\cite{lamport1978time}: event $a$
happened before event $b$ if they occur in that order in the same process, or
if $a$ is the sending of a message that $b$ receives, or transitively through
a chain of such steps (Figure~\ref{fig:hb}). Happened-before is a partial
order (some event pairs are simply concurrent), and it is the order
that determines whether one event could have influenced another: in
Figure~\ref{fig:hb}, $a$ happened before $d$ through the chain
$a \to c \to e \to d$, while $b$ and $e$ are concurrent: no causal path
connects them, even though $b$ occurs later in wall-clock time.

\begin{figure}[t]
  \centering
  % Space-time diagram: happened-before, causal chains, concurrency.
\begin{tikzpicture}[>=Stealth, thick, font=\small]
  % process timelines
  \foreach \y/\p in {0/{$P_3$}, 1.5/{$P_2$}, 3.0/{$P_1$}} {%
    \draw[->, semithick] (0,\y) -- (10.4,\y);
    \node[left=3pt] at (0,\y) {\p};
  }
  \node[below right] at (9.4,-0.1) {\footnotesize wall-clock time};
  % events
  \filldraw (1.2,3.0) circle (2.2pt) node[above=2pt] {$a$};
  \filldraw (5.2,3.0) circle (2.2pt) node[above=2pt] {$b$};
  \filldraw (2.6,1.5) circle (2.2pt) node[above=2pt] {$c$};
  \filldraw (6.6,1.5) circle (2.2pt) node[above=2pt] {$d$};
  \filldraw (4.2,0)   circle (2.2pt) node[below=2pt] {$e$};
  \filldraw (8.6,0)   circle (2.2pt) node[below=2pt] {$f$};
  % messages
  \draw[->] (1.2,3.0) -- node[midway, left=2pt]  {\footnotesize $m_1$} (2.6,1.5);
  \draw[->] (2.6,1.5) -- node[midway, left=2pt]  {\footnotesize $m_2$} (4.2,0);
  \draw[->] (4.2,0)   -- node[midway, right=3pt] {\footnotesize $m_3$} (6.6,1.5);
\end{tikzpicture}

  \caption{A space-time diagram of three processes exchanging messages
  $m_1$--$m_3$. Causality flows along process lines and messages:
  $a \rightarrow c \rightarrow e \rightarrow d$ is a causal chain, so $a$
  \emph{happened before} $d$; $b$ and $e$ are \emph{concurrent} (neither
  could have influenced the other), despite $b$ being later on the wall
  clock.}
  \label{fig:hb}
\end{figure}

Logical and vector
clocks~\cite{lamport1978time,fidge1988timestamps,mattern1989virtual} make
this order computable
from local information; alternatively, when all events carry timestamps from
reasonably synchronized physical clocks and land in one ordered store, causal
chains follow from timestamps plus application-level correlation identifiers.

The corollary this thesis leans on is that the
\emph{event log}, an append-only sequence of timestamped, attributed
events, is the canonical representation of distributed execution. A log
captures exactly what each process observed and when; every richer structure
(a request trace, a communication graph, a health summary) is a fold over it.

\emph{Used in this thesis:} multi-agent executions are treated as partially
ordered event logs; causal structure is recovered from physical-time ordering
plus explicit correlation identifiers carried by the instrumentation
(Chapters~\ref{chap:design} and~\ref{chap:analytics}).

\subsection{Delivery Semantics}
\label{sec:bg-delivery}

Any pipeline that moves events from producers to a store must define what
happens when a component fails mid-transfer. Three contracts are
standard~\cite{kleppmann2017ddia}.
\emph{At-most-once} delivery never retries, so failures lose events,
which is unacceptable for an audit trail. \emph{At-least-once} delivery retries until
acknowledged, so failures duplicate events. \emph{Exactly-once} delivery is
the ideal, unachievable in general between separately failing
components; real systems approximate it as \emph{effectively exactly-once}:
at-least-once transport paired with deduplication at the consumer, using
idempotent writes or monotonic sequence numbers (``high-water marks'') so
that redelivered events are recognized and dropped.

A related design decision is where events wait when a downstream component is
slow or dead. Buffering writes in a local, durable \emph{spool} decouples the
producer from downstream availability: the producer's fast path is a local
append, and a separate drain process moves data onward at its own pace,
the same write-ahead pattern that underlies journaling file systems and
message brokers.

\emph{Used in this thesis:} capture is at-least-once from a durable local
spool; ingest into the log is made effectively exactly-once by per-session
high-water marks, with consumer state rebuilt from the log itself after a
crash (\S\ref{sec:design-write}).

\section{Distributed Tracing}
\label{sec:bg-tracing}

\subsection{Lineage: From Dapper to Pixie}
\label{sec:bg-tracing-lineage}

Distributed tracing is the discipline of reconstructing, from per-process
records, the cross-machine causal story of a single request. Google's
Dapper~\cite{sigelman2010dapper} established the modern vocabulary: a
\emph{trace} is a tree of \emph{spans}, each span a timed operation carrying
the trace's identifier, its parent's identifier, and annotations. Two Dapper
design tenets recur throughout this thesis. First, instrumentation must be
low-overhead and effectively invisible to application developers, or it will
be turned off or bypassed; Dapper achieved this by instrumenting shared
infrastructure (the RPC library) rather than applications. Second, complete
capture at scale is considered infeasible, so Dapper samples aggressively
(uniformly, in its default scheme).

Subsequent systems relaxed different assumptions. The Mystery
Machine~\cite{chow2014mystery} showed that when explicit propagation is
unavailable, the causal model itself can be \emph{inferred}: from large
numbers of logged executions, it hypothesizes happened-before relationships
and prunes those contradicted by observed orderings. Pivot
Tracing~\cite{mace2015pivot} made tracing \emph{dynamic}: operators install
queries at runtime that propagate context along the execution's causal path
(the ``baggage'' abstraction), answering questions that were never
pre-instrumented. Canopy~\cite{kaldor2017canopy} scaled the pipeline
end-to-end at Facebook, emphasizing that raw traces are not what engineers
consume: traces are transformed into feature-indexed datasets that back
purpose-built views. More recently, systems such as Pixie~\cite{pixie} pushed
capture into the platform itself via eBPF, trading application-level
semantics for zero-touch deployment.

\subsection{Critical-Path Analysis and Self-Time}
\label{sec:bg-critpath}

A trace answers ``what happened''; the follow-up question is ``why was it
slow.'' In a tree of overlapping timed operations, total latency is not the
sum of span durations: children execute inside their parents, and siblings
overlap. The \emph{critical path}, an idea imported into parallel-program
performance analysis by Yang and Miller~\cite{yang1988critpath}, is the
chain of operations from root to
leaf whose durations determine end-to-end latency: shortening any
operation off the critical path cannot speed up the request. Within the
critical path, the correct measure of each hop's contribution is its
\emph{self-time} (exclusive time): wall-clock time spent in the operation
itself, excluding intervals covered by its children. Attributing latency by
inclusive duration double-counts every level of the tree; self-time
attribution names the hop actually responsible. Figure~\ref{fig:critpath}
illustrates both ideas on a small trace: shortening $B$ or $C_1$ cannot
speed up the request, and although $C$'s inclusive duration is 720\,ms,
only 320\,ms of it is $C$'s own.
The Mystery Machine used this analysis to find optimization
opportunities at Facebook~\cite{chow2014mystery}.

\begin{figure}[t]
  \centering
  % Trace waterfall: span tree, critical path (grey), self-time.
\begin{tikzpicture}[font=\footnotesize, thick, >=Stealth]
  % time axis
  \draw[->, semithick] (0,-0.6) -- (12.6,-0.6);
  \node[below] at (11.8,-0.65) {time (ms)};
  \foreach \x/\t in {0/0, 3/300, 6/600, 9/900, 12/1200} {
    \draw (\x,-0.52) -- (\x,-0.68);
    \node[below] at (\x,-0.7) {\tiny \t};
  }
  % depth 0: root span A (critical)
  \draw[very thick, fill=gray!25] (0,3.2) rectangle (12,3.85);
  \node at (6,3.52) {A \; (orchestrator) \; 1200\,ms};
  % depth 1: B (off-path), C (critical)
  \draw[fill=white] (0.4,2.3) rectangle (4.4,2.95);
  \node at (2.4,2.62) {B \; 400\,ms};
  \draw[very thick, fill=gray!25] (4.6,2.3) rectangle (11.8,2.95);
  \node at (8.2,2.62) {C \; 720\,ms};
  % depth 2: C1 (off-path), C2 (critical)
  \draw[fill=white] (5.0,1.4) rectangle (7.0,2.05);
  \node at (6.0,1.72) {C$_1$ \; 200\,ms};
  \draw[very thick, fill=gray!25] (7.5,1.4) rectangle (11.5,2.05);
  \node at (9.5,1.72) {C$_2$ \; 400\,ms};
  % self-time of C: its interval not covered by its slowest child C2
  \draw[decorate, decoration={brace, mirror, raise=3pt}] (4.6,1.4) -- (7.5,1.4);
  \node[below=9pt] at (6.05,1.4)
    {self-time of C $= 720 - 400 = 320$\,ms};
  % legend
  \draw[very thick, fill=gray!25] (0.0,0.1) rectangle (0.5,0.4);
  \node[right] at (0.55,0.25) {critical path};
  \draw[fill=white] (3.2,0.1) rectangle (3.7,0.4);
  \node[right] at (3.75,0.25) {off the critical path};
\end{tikzpicture}

  \caption{A trace rendered as a waterfall. Grey spans form the critical
  path $A \rightarrow C \rightarrow C_2$, the chain that determines
  end-to-end latency; white spans ($B$, $C_1$) overlap it and cannot
  shorten the request. Each hop is charged its \emph{self-time}: $C$'s
  720\,ms minus the 400\,ms covered by its slowest child leaves 320\,ms
  of $C$'s own.}
  \label{fig:critpath}
\end{figure}

\emph{Used in this thesis:} the tracing view of Chapter~\ref{chap:analytics}
reconstructs cross-host call trees from inter-agent events (using explicit
parent correlation when present and Mystery-Machine-style inference from
session and time nesting when not), computes the critical path, and names
the bottleneck hop by self-time.

\subsection{Sampling and Its Limits at Fleet Scale}
\label{sec:bg-sampling}

Dapper's sampling assumption does not transfer cleanly to the LLM-agent
setting. Head-based sampling (deciding at trace start) keeps coherent
traces but misses rare failures by construction. Tail-based sampling
(deciding at trace end) must buffer every trace until its fate is known.
At fleet scale that reintroduces the data-volume problem sampling was
meant to solve, and it requires reassembling spans scattered across many
hosts before the keep/drop decision can be made.
LLM-agent workloads change the calculus in two ways.
First, event rates are modest: an agent turn takes seconds and costs real
money, so a fleet of hundreds of agents emits orders of magnitude fewer
events than a microservice fleet emits RPCs\footnote{The arithmetic: a
100-agent fleet at one turn per agent every few seconds produces tens of
events per second; RPC tracing operates against request rates many orders
of magnitude higher, which is why Dapper samples at all.}; complete
capture is affordable. Second, the value of completeness is higher: each event is
semantically rich (a full LLM interaction), failures are the rare
events sampling discards, and post-hoc analyses such as cost accounting
require totals, not estimates.

\emph{Used in this thesis:} the system captures completely and treats
storage, not sampling, as the scaling lever, pushing the volume problem
down to the log substrate and the read-side rollups
(\S\ref{sec:bg-log},~Chapter~\ref{chap:analytics}).

\section{The Distributed Shared Log Abstraction}
\label{sec:bg-log}

The systems of \S\ref{sec:bg-tracing} all separate capture from storage: spans
are generated in-process, then shipped to a backend whose design is
incidental to the tracing model. A complementary line of work treats the
\emph{log itself} as the system's central abstraction. Kreps's
formulation~\cite{kreps2013log} argues that an append-only, totally ordered
log is the natural integration point for distributed data systems: producers
append, consumers read at their own pace, and the log's order \emph{is} the
system's notion of time (Figure~\ref{fig:sharedlog}).
Kafka~\cite{kreps2011kafka} industrialized this as
partitioned commit logs (replication followed in later versions); CORFU~\cite{balakrishnan2012corfu}
showed that a shared log can be a first-class distributed primitive:
clients append to a global address space striped over a cluster, and the
shared log's total order becomes the foundation for replicated state
machines and transactional systems built above it.

\begin{figure}[t]
  \centering
  % The shared log: producers append at the tail, decoupled readers at offsets.
\begin{tikzpicture}[font=\footnotesize, thick, >=Stealth]
  % ordered cells
  \foreach \i in {0,...,9} {
    \draw ({\i*1.1},0) rectangle ({\i*1.1+1.1},0.8);
    \node at ({\i*1.1+0.55},0.4) {$e_{\i}$};
  }
  \draw[dashed] (11.0,0) rectangle (12.1,0.8);
  % order annotation
  \draw[->, semithick] (0,1.25) -- (12.1,1.25);
  \node[above] at (6.05,1.28) {append-only; position in the log $=$ order of history};
  % producers appending at the tail
  \node[draw, fill=white] (pa) at (8.6,2.6)  {producer $A$};
  \node[draw, fill=white] (pb) at (11.6,2.6) {producer $B$};
  \draw[->] (pa.south) .. controls (10.3,1.9) .. (10.45,0.85);
  \draw[->] (pb.south) .. controls (11.6,1.9) .. (11.55,0.85);
  \node at (10.1,2.6) {append};
  % decoupled readers at their own offsets
  \node[draw, fill=white] (r1) at (2.75,-1.6) {reader 1 (offset 2)};
  \node[draw, fill=white] (r2) at (7.6,-1.6)  {reader 2 (offset 6)};
  \draw[->] (r1.north) -- (2.75,-0.05);
  \draw[->] (r2.north) -- (7.15,-0.05);
\end{tikzpicture}

  \caption{The shared-log abstraction. Producers append at the tail;
  an event's position in the log is its permanent place in history.
  Readers consume independently at their own offsets, imposing no
  back-pressure on producers or on each other.}
  \label{fig:sharedlog}
\end{figure}

Three properties of this abstraction matter for observability. First,
\emph{node-local appends}: producers write to a nearby log segment, keeping
the fast path off the network's critical path. Second, \emph{decoupled
readers}: consumers impose no back-pressure on producers; an expensive
analysis reads the same log as a cheap dashboard without perturbing capture.
Third, \emph{durability with order}: once appended, an event's position,
its place in history, is permanent, the property an audit trail of
agent behavior needs.

The abstraction also has a characteristic trade-off: log stores are
\emph{write-optimized}. Appends are fast because data lands in
ingest-side buffers that are sealed, ordered, and moved to
read-optimized storage asynchronously; a reader wanting the most recent
events must either wait out this pipeline or read from a separate,
freshness-oriented structure. Stream-processing architecture has long
formalized the resulting design: the \emph{lambda architecture} pairs a
durable, complete, high-latency batch path with a fast, approximate speed
layer over the same event stream~\cite{marz2015big}. This trade-off, in
ChronoLog's concrete form, drives a central design decision of this thesis
(\S\ref{sec:bg-chronolog-design}, Chapter~\ref{chap:design}).

\emph{Used in this thesis:} the shared log is the single substrate
(capture appends to it, every view reads from it), and its
write-optimized/read-lagging nature is handled by a lambda-style hot/cold
split rather than by abandoning the abstraction.

\section{ChronoLog}
\label{sec:bg-chronolog}

ChronoLog~\cite{kougkas2020chronolog} is a distributed, tiered shared-log
store developed for HPC environments, and the concrete substrate of this
thesis. The paper presents the design; the component names, tier formats,
and configuration described below reflect the deployed version used in
this work. It differs from the logs of \S\ref{sec:bg-log} in one decision that
makes it particularly suited to this thesis's workload: events are ordered by
\emph{physical time} at capture, rather than by a sequencer-assigned logical
position, so the log's global order directly reflects when things happened
across the cluster.

\subsection{Architecture: Visor, Keepers, Grapher--Player}
\label{sec:bg-chronolog-arch}

A ChronoLog deployment (Figure~\ref{fig:chronoarch}) comprises a
coordinating \emph{visor}, a recording group of \emph{keepers}, and an
archival pipeline of \emph{grapher} and \emph{player} processes.

\begin{figure}[t]
  \centering
  % ChronoLog: visor / keeper-per-node / grapher drain / archive.
\begin{tikzpicture}[font=\footnotesize, thick, >=Stealth,
    comp/.style={draw, fill=white, minimum height=0.7cm, minimum width=1.7cm, align=center}]
  % three cluster nodes, each with local agents appending to a local keeper
  \foreach \i/\y in {1/4.4, 2/2.2, 3/0} {
    \node[comp] (ag\i) at (0,\y)   {agents};
    \node[comp] (kp\i) at (3.2,\y) {keeper \i};
    \draw[->] (ag\i) -- node[above] {\tiny local append} (kp\i);
    \node[draw, dashed, fit=(ag\i)(kp\i), inner sep=6pt] (nb\i) {};
    \node[anchor=south west, inner sep=2pt] at (nb\i.north west) {\tiny compute node \i};
  }
  % visor: control plane
  \node[comp] (visor) at (7.0,4.9) {visor};
  \node[right=2pt of visor, align=left] {\tiny control plane:\\ \tiny connect, acquire,\\ \tiny discover keepers};
  \foreach \i in {1,2,3} { \draw[dashed, semithick] (visor.west) -- (kp\i.north east); }
  % grapher drain -> archive
  \node[comp] (gr) at (7.0,2.2)  {grapher};
  \node[comp, minimum width=2.2cm] (ar) at (10.6,2.2) {archive\\ (CSV/HDF5)};
  \foreach \i in {1,2,3} { \draw[->] (kp\i.east) -- (gr.west); }
  \draw[->] (gr) -- node[above] {\tiny merge, order} (ar);
  \node[below=2pt, align=center] at (gr.south)
    {\tiny asynchronous drain:\\ \tiny acceptance window $\approx$ 180\,s};
  % readers
  \node[comp] (rd) at (10.6,0) {readers\\ (replay)};
  \draw[->] (ar) -- (rd);
\end{tikzpicture}

  \caption{ChronoLog's architecture. A keeper on every compute node takes
  node-local appends into memory (the ingest tier); the visor coordinates
  clients; the grapher asynchronously drains, merges, and time-orders
  sealed chunks into a read-optimized archive. The drain imposes an
  acceptance window during which fresh events are not yet readable
  (20--30\,s measured end-to-end here, with a 180\,s window configured
  at the grapher). This is the property the read path of
  Chapter~\ref{chap:design} is designed around.}
  \label{fig:chronoarch}
\end{figure}

The visor is the control plane: clients connect to
it to create and acquire logs and to discover the recording group. The
keepers are the ingest plane: one keeper runs on every node of the
deployment, and clients append events to a node-local keeper, which buffers
recent history in memory. Asynchronously, the grapher drains sealed chunks
from the keepers, merges them into time order, and writes them to a
read-optimized archive (CSV/HDF5 tiers), which the player serves back to
readers. The design mirrors the tiered storage hierarchies of HPC I/O
systems: a fast, distributed, memory-resident ingest tier in front of a
durable, ordered archival tier.

Two consequences of this architecture are load-bearing for this thesis. The
keeper-per-node layout means capture is \emph{physically co-located} with the
workload: an agent's events are appended on the machine where the agent runs,
and the mapping from events to nodes (the coupling of semantics to
topology that Chapter~\ref{chap:intro} identified as the gap) is native to
the substrate rather than reconstructed. And the drain pipeline means the
archive is complete and ordered but \emph{late} (\S\ref{sec:bg-chronolog-design}).

\subsection{Data Model: Chronicles, Stories, Events}
\label{sec:bg-chronolog-model}

ChronoLog organizes data hierarchically. A \emph{chronicle} is a named log
namespace; within it, a \emph{story} is an independently written and read
event sequence; an \emph{event} is a timestamped payload appended to a story.
Clients acquire a story handle to append or replay, and stories provide the
unit of isolation: writers to different stories do not contend, and a reader
replays exactly one story's history in time order. The model maps naturally
onto observability data: chronicles partition by data kind, stories by
session or entity, and events carry the individual observations
(\S\ref{sec:design-mapping} details this thesis's schema).

\subsection{Design Point: Ingest-Optimized, Delayed Reads}
\label{sec:bg-chronolog-design}

ChronoLog's design point is high-throughput, low-latency \emph{ingest}:
appends land in the node-local keeper's memory and return. The price is paid
on the read side, in two forms that this thesis measures and designs around.
First, the keeper$\to$grapher drain imposes an \emph{acceptance window}
during which freshly appended events are held in ingest buffers and are
not yet readable from the archive: in the deployments used here the
keeper seals 10\,s chunks under a 15\,s acceptance window, for a
measured end-to-end availability of 20--30\,s
(\S\ref{sec:eval-freshness}); the grapher and player are configured
with a longer, 180\,s window. Second, a live replay can be not merely
stale but \emph{blocking}: a read that the player cannot answer holds
the client (and, in a Python client holding the global interpreter
lock, the entire calling process) for the full ${\sim}180$\,s query
timeout. A reader architecture built naively on such a
substrate freezes exactly when it is most interesting: while the system is
live.

This is not a defect but the standard write-optimized trade-off of
\S\ref{sec:bg-log} in concrete form, and it dictates the read-side
architecture of Chapter~\ref{chap:design}: ChronoLog serves as the durable,
complete cold path, while real-time views are fed by a hot path populated at
ingest. That is the lambda split, with the acceptance window as the boundary
between the layers.

\section{LLM Agents and Multi-Agent Systems}
\label{sec:bg-agents}

\subsection{The Agent Loop, Tool Use, and Delegation}
\label{sec:bg-agent-loop}

An LLM \emph{agent} is a program that couples a language model to an
environment through a loop (Figure~\ref{fig:agentloop}a): the model is
prompted with a goal and a transcript of prior steps; it emits either a
final answer or an \emph{action}, a call to an external tool; the tool's
result is appended to the transcript; and the loop repeats. The pattern was
crystallized by ReAct~\cite{yao2023react}, which interleaves free-form
reasoning with structured actions. The unit of execution is the
\emph{turn}: one model invocation, with its prompt, response, latency, and
token consumption.

\begin{figure}[t]
  \centering
  % (a) the single-agent loop; (b) multi-agent delegation across nodes.
\begin{tikzpicture}[font=\footnotesize, thick, >=Stealth,
    comp/.style={draw, fill=white, minimum height=0.65cm, minimum width=1.5cm, align=center},
    agent/.style={draw, fill=white, circle, minimum size=0.8cm}]
  % ---- panel (a): the agent loop ----
  \node[comp] (ctx)  at (0,0)     {context};
  \node[comp] (llm)  at (3.1,1.0) {LLM};
  \node[comp] (tool) at (3.1,-1.0){tool};
  \draw[->] (ctx.north) to[bend left=25] node[above left=-2pt] {\tiny prompt} (llm.west);
  \draw[->] (llm.south) to[bend left=25] node[right=2pt] {\tiny action} (tool.north);
  \draw[->] (tool.west) to[bend left=25] node[below left=-2pt] {\tiny observation} (ctx.south);
  \draw[->] (llm.east) -- +(1.0,0) node[right] {\tiny answer};
  \node at (1.9,-2.2) {(a) the agent loop: one \emph{turn} per LLM call};
  % ---- panel (b): delegation across nodes ----
  \node[agent] (A) at (7.6,0.0)   {$A$};
  \node[agent] (B) at (11.0,0.9)  {$B$};
  \node[agent] (C) at (11.0,-0.9) {$C$};
  \node[draw, dashed, fit=(A), inner sep=8pt] (n1) {};
  \node[anchor=south west, inner sep=2pt] at (n1.north west) {\tiny node 1};
  \node[draw, dashed, fit=(B)(C), inner sep=8pt] (n2) {};
  \node[anchor=south west, inner sep=2pt] at (n2.north west) {\tiny node 2};
  \draw[->] (A) -- node[above, sloped] {\tiny start} (B);
  \draw[->, dashed] (B) to[bend right=18]
    node[above, sloped] {\tiny done (status, latency)} (A);
  \draw[->] (B) -- node[right] {\tiny start} (C);
  \draw[->, dashed] (C) to[bend left=30] (B);
  \node at (9.7,-2.2) {(b) delegation: agents as tools, calls across nodes};
\end{tikzpicture}

  \caption{(a) The agent loop: each cycle through the LLM is one
  \emph{turn}, appending the tool's observation to the agent's context.
  (b) Multi-agent delegation: once agents are tools, agent $A$ calls agent
  $B$ (possibly on another node), which calls $C$ in turn, each call a
  \texttt{start} with a matching \texttt{done} (dashed) carrying status and
  latency. The result is a partially ordered event log with delegations as
  its cross-process edges.}
  \label{fig:agentloop}
\end{figure}

Tool interfaces are standardizing around protocols such as the Model Context
Protocol~\cite{anthropic2024mcp}, in which tools are typed, discoverable
endpoints. Once tools are first-class, \emph{agents themselves can be tools}:
one agent's toolset includes a call-remote-agent operation that opens a
sub-conversation with another agent, possibly on another machine. This is
how multi-agent systems form~\cite{wu2023autogen}: a
scenario is a set of agents, each a loop over its own context, connected by
delegation edges. From a distributed-systems standpoint
(\S\ref{sec:bg-time}), a scenario execution is a set of per-agent event
sequences joined by inter-agent messages: precisely a partially ordered
event log, with delegation calls as the cross-process edges.

\subsection{Observable Signals}
\label{sec:bg-agent-signals}

The observable surface of an agent fleet has two planes. The
\emph{intra-agent} plane is the sequence of LLM turns: prompt and response
content, model identity, latency, token counts, and (unusual among
workloads) direct monetary cost per event, since providers bill per token.
It also includes the agent's \emph{working memory}: the context assembled for
each turn, which grows as results accumulate and shrinks under eviction and
summarization when the model's context window fills. Memory dynamics are an
operational signal in their own right: an agent nearing its context limit
degrades before it fails. The \emph{inter-agent} plane is the set of
delegation calls: which agent called which, when the call started and
finished, whether it succeeded, and how calls nest. This is the plane from
which scenario structure, communication topology, and cross-host traces are
reconstructed.

\subsection{Failure Modes of LLM Agent Fleets}
\label{sec:bg-agent-failures}

Agent fleets fail in ways that combine both planes. Empirical taxonomies of
multi-agent LLM failures~\cite{cemri2025mast} find the dominant categories
to be specification gaps, inter-agent misalignment, and missing task
verification, all of which surface to an operator as the concrete symptoms
below. On the intra-agent plane:
provider errors (HTTP 4xx/5xx, rate limits), latency outliers against a
model's typical distribution, and \emph{retry storms}, in which an agent
re-issues a failing prompt in a loop, a classic amplification pathology
(the retry-driven feedback loops of metastable
failure~\cite{bronson2021metastable}), here with a direct monetary cost
attached to every retry. Agents may
also \emph{backtrack}, restoring earlier context after an unproductive path;
frequent recoveries signal a struggling agent. On the inter-agent plane:
failed delegations, and, particularly damaging, \emph{hung} calls
(termed \emph{orphan calls} from Chapter~\ref{chap:analytics} on), in
which a delegation starts but never completes, silently stalling the
delegating agent; in a per-edge view nothing is visibly wrong, since the
signature is an absence. At fleet scale these interact with infrastructure
pathology: a degraded, fail-slow node~\cite{gunawi2018failslow} drags down
every agent scheduled on it, producing a symptom that is scattered across
scenarios yet has a single physical cause. Localizing such failures is
the cross-layer question neither tooling camp of
\S\ref{sec:bg-sota} can answer alone.

\section{State of the Art and the Gap}
\label{sec:bg-sota}

\subsection{LLM and Agent Observability}
\label{sec:bg-sota-llm}

The first camp instruments the application layer.
LangSmith~\cite{langsmith} (from the LangChain ecosystem),
Langfuse~\cite{langfuse}, and Arize Phoenix~\cite{phoenix} capture LLM calls
and tool invocations as hierarchical runs or OpenTelemetry-style spans,
adding LLM-specific value: token and cost accounting, prompt management,
dataset curation, and LLM-as-judge evaluation. AgentOps~\cite{agentops}
targets agent frameworks specifically, tracking sessions, tool usage, and
multi-agent hand-offs. General application-performance-monitoring (APM)
vendors have added LLM-observability products~\cite{datadogllm}, and the
OpenTelemetry project is standardizing generative-AI semantic
conventions~\cite{otelgenai}, signalling consolidation of this camp around
the span model.

The limitation is structural, not incidental. The span model these tools
inherit was designed for request-scoped microservice tracing:
``distributed'' means span correlation across services in one
application's request path; ``multi-agent'' means multiple agent
abstractions within one process or one orchestrator. Host identity, where
it appears at all, is opaque metadata: a tag an operator can attach and
read back, not a first-class dimension the tools can group by, compare
across, or alert over. Population-level questions (compare nodes,
find the sick one, correlate agent failures with machine placement)
therefore cannot be phrased. The nearest exceptions prove the shape of the gap
rather than filling it: an APM platform such as Datadog can correlate an
LLM trace with the host metrics of \emph{that request's} infrastructure,
but offers no fleet-of-agents semantics over an HPC allocation: no
batch-scheduler awareness, no population view keyed by agent behavior.
The OTel generative-AI conventions standardize \emph{turn} semantics,
while the host resource attributes an SDK may attach come with no defined
fleet-level analysis. In both cases the data model cannot express the join
the operator questions of \S\ref{sec:intro-context} require: from a
fleet-wide symptom to a specific agent's conversation on a named machine.

\subsection{HPC and Fleet Monitoring}
\label{sec:bg-sota-hpc}

The second camp instruments the physical layer, and does so at scale.
Prometheus~\cite{prometheus}, typically paired with Grafana dashboards, is
a de facto standard for metrics-based monitoring, routinely deployed over
SLURM clusters via node exporters and job accounting; HPC-native systems such as
LDMS~\cite{agelastos2014ldms} are engineered for continuous, low-overhead
collection of node metrics across very large machines. These systems answer
per-node and per-job questions (utilization, saturation, hardware
health, stragglers) well, and their visual idioms (heatmaps, per-node grids) are
built for hundreds of nodes.

Their blindness is the mirror image of the first camp's. The data model is
metric time series named by host and counter; there is no application
semantics at all. That a CPU-saturated node is running an LLM agent whose
retry storm is the actual pathology (and that the same storm is visible as
cost and latency in some scenario's transcript) is invisible: metrics
cannot say \emph{which conversation} is responsible, and offer no path from
a red node to the agent behavior causing it.

\subsection{The Empty Intersection}
\label{sec:bg-sota-gap}

Figure~\ref{fig:quadrant} summarizes the landscape along its two axes:
richness of agent semantics, and awareness of physical fleet topology. Camp
1 occupies the semantics-rich, topology-blind quadrant; Camp 2 the
topology-aware, semantics-blind one. The quadrant combining both (agent
semantics coupled to cluster topology, at fleet scale) is, among the
tools surveyed here, unoccupied. It
is where the operator questions of \S\ref{sec:intro-context} live: they
require pivoting \emph{between} layers, from a fleet-level symptom to a
single agent's conversation and back.

\begin{figure}[t]
  \centering
  % The capability quadrant: the empty intersection is this thesis.
\begin{tikzpicture}[font=\footnotesize, thick, >=Stealth]
  % axes
  \draw[->] (-0.35,-0.35) -- (11.1,-0.35);
  \node[below] at (5.3,-0.45) {awareness of physical fleet topology $\longrightarrow$};
  \draw[->] (-0.35,-0.35) -- (-0.35,6.7);
  \node[rotate=90, above] at (-0.5,3.1) {richness of agent semantics $\longrightarrow$};
  % bottom-left: neither
  \draw (0,0) rectangle (5.2,3.0);
  \node[align=center] at (2.6,1.5) {---};
  % bottom-right: Camp 2
  \draw (5.5,0) rectangle (10.7,3.0);
  \node[align=center] at (8.1,1.5)
    {\textbf{Camp 2: fleet monitoring}\\[2pt]
     Prometheus/Grafana on SLURM\\ LDMS, HPC monitoring\\[3pt]
     \emph{node health at scale,}\\ \emph{no agent semantics}};
  % top-left: Camp 1
  \draw (0,3.3) rectangle (5.2,6.3);
  \node[align=center] at (2.6,4.8)
    {\textbf{Camp 1: LLM observability}\\[2pt]
     LangSmith, Langfuse, Phoenix,\\ AgentOps, Datadog LLM, OTel GenAI\\[3pt]
     \emph{rich agent semantics,}\\ \emph{no physical topology}};
  % top-right: the gap / this thesis
  \draw[very thick, fill=gray!20] (5.5,3.3) rectangle (10.7,6.3);
  \node[align=center] at (8.1,4.8)
    {\textbf{this thesis}\\[2pt]
     agent semantics coupled to\\ cluster topology, at fleet scale\\[3pt]
     \emph{previously unoccupied}};
\end{tikzpicture}

  \caption{The observability landscape for multi-agent LLM systems on
  clusters, along its two capability axes. Camp 1 and Camp 2 each occupy
  one off-diagonal quadrant; their intersection (agent semantics coupled
  to physical fleet topology) is unoccupied among the surveyed tools,
  and is the quadrant this thesis targets.}
  \label{fig:quadrant}
\end{figure}

Filling that quadrant is not a matter of running both camps side by side:
two disjoint systems with disjoint data models provide no join between an
agent event and a node, which is precisely the pivot the questions require.
The evaluation returns this structural argument as a measurement
(\S\ref{sec:eval-casestudy}): faced with an injected cross-layer fault on
a live cluster, an operator given the raw captured files (all the
signals both camps would collect, join included) could not produce a
single localization answer. Through the joined views, that operator
completed the full symptom-to-conversation pivot in 72 seconds.
The same experiment also separates the gap's two halves: sufficiency is a
property of the \emph{data model} (an expert reader extracted every
answer from the log files alone, because host attribution and causal
correlation are in the events themselves; \S\ref{sec:eval-casestudy}
gives the details), while accessibility is a property of the
\emph{interface} built on that model. The camps fail at the first half,
which no interface can repair: a view over data that lacks the join
cannot invent it.
The gap therefore calls for a single substrate in which agent-semantic
events are natively attributed to physical nodes: what the
shared-log approach of \S\ref{sec:bg-log}, instantiated on ChronoLog's
keeper-per-node architecture (\S\ref{sec:bg-chronolog}), provides by
construction. Designing the capture, read, and analysis layers that turn
that substrate into a working observability system is the subject of the
remainder of this thesis.
